\section{Discussion}

The analysis indicates that molecular property prediction benchmarks should be interpreted as a joint product of the model, the dataset, and the train/test split. Random splitting, ordinary scaffold splitting, and target-balanced scaffold splitting do not provide interchangeable test sets. They produce different chemical similarity patterns, target distributions, scaffold compositions, performance gaps, and, in some cases, model rankings. A benchmark score is therefore incomplete without information about the split that produced it.

Random-split performance should be interpreted cautiously in molecular property prediction. In this audit, random splits generally produced higher maximum test-to-train Tanimoto similarity than scaffold-based splits. Random splits remain useful for measuring interpolation-like performance under chemically similar train/test distributions, but they should not be interpreted as evidence of scaffold-level chemical generalization unless chemical similarity is also reported.

Ordinary scaffold splitting is useful, but it is not self-explanatory. It removes train/test scaffold overlap, which is valuable for evaluating scaffold-level separation. At the same time, it can induce target shift and scaffold concentration. ESOL and FreeSolv illustrate this problem clearly: the ordinary scaffold split holds out very few scaffold groups and creates large target gaps. In such cases, a performance drop under scaffold splitting may reflect a mixture of scaffold generalization, target distribution shift, and small-test-set instability. Reporting scaffold statistics and target gaps makes this ambiguity visible.

The target-balanced scaffold split helps clarify this ambiguity. When gaps shrink after target balancing, as observed for several ESOL tree-based results and some classification models, part of the ordinary scaffold gap is plausibly associated with target shift. When gaps remain or increase, as observed in HIV, the split may be probing a different difficulty regime. Target-balanced scaffold splitting should therefore be interpreted as a diagnostic counterfactual rather than as a universal benchmark replacement.

For molecular property prediction benchmarks, model scores should be accompanied by split diagnostics. At minimum, reports should include the split protocol, target distribution diagnostics, scaffold statistics, test-to-train chemical similarity, and model-ranking stability. These additions do not require more complex models, but they make model comparisons easier to interpret and reduce the risk of attributing split-induced artifacts to model performance.

\section{Limitations}

This study has several limitations. First, the experiments use fingerprint-based classical baselines rather than graph neural networks or transformer-based molecular encoders. This was intentional: lightweight baselines isolate split effects and improve reproducibility. Nevertheless, future work should test whether the same diagnostic patterns hold for deep molecular representations. Second, the target-balanced scaffold split is heuristic. It is used here as a diagnostic counterfactual, not as a new universal split standard. Third, the benchmark includes six public datasets, which provide a focused audit across classification and regression tasks but do not cover all AIDD scenarios. Fourth, no prospective or time-based validation is included. Finally, small datasets such as FreeSolv may be structurally constrained by scaffold concentration, making it difficult to construct a split that is simultaneously scaffold-separated, target-balanced, and statistically stable.

\section{Conclusion}

Molecular property prediction benchmarks should treat data splitting as an object of diagnosis. Random splits can produce chemically closer test sets and optimistic performance estimates. Ordinary scaffold splits remove scaffold overlap but can introduce target shift and scaffold concentration. Target-balanced scaffold splits help identify whether scaffold gaps persist after reducing target shift, but their effects remain dataset- and model-dependent. A split-diagnostic audit that reports performance together with target distribution shift, scaffold statistics, chemical similarity, outlier cases, and model-ranking stability provides a more transparent basis for evaluating generalization in AI-driven drug discovery.
